\newpage
\phantomsection
\label{prf:rolles-theorem}

\begin{remark*}[Return]
\hyperref[thm:rolles-theorem]{Return to Theorem}
\end{remark*}

\begin{theorem*}[Rolle's Theorem]
Let $a, b \in \mathbb{R}$ with $a < b$. Suppose that $f : [a,b] \to \mathbb{R}$ is continuous on $[a,b]$, differentiable on $(a,b)$, and satisfies $f(a) = f(b)$. Then there exists $c \in (a,b)$ such that $f'(c) = 0$.
\end{theorem*}

\begin{proof}
\textbf{Professional Standard Proof.}~\\
By the Extreme Value Theorem, $f$ attains a maximum $M$ and a minimum $m$ on $[a,b]$. If $M = m$ then $f$ is constant on $[a,b]$, so $f' \equiv 0$ on $(a,b)$ and any $c \in (a,b)$ serves. Otherwise $M > m$, so at least one of $M, m$ differs from $f(a) = f(b)$ and is therefore attained at some interior point $c \in (a,b)$. Since $c$ is a relative extremum and $f$ is differentiable at $c$, the Interior Extremum Theorem gives $f'(c) = 0$.
\end{proof}

\begin{proof}
\textbf{Detailed Learning Proof.}~\\

\textbf{Step 1.} Apply the Extreme Value Theorem. Since $f$ is continuous on the closed bounded interval $[a,b]$, the Extreme Value Theorem guarantees that $f$ attains both its maximum value $M = \max_{[a,b]} f$ and its minimum value $m = \min_{[a,b]} f$ on $[a,b]$.

\textbf{Step 2.} Handle the degenerate case. If $M = m$, then $f$ is constant on $[a,b]$ --- every value of $f$ equals $M = m$. A constant function has derivative zero everywhere it is differentiable. Therefore $f'(c) = 0$ for every $c \in (a,b)$, and the theorem holds.

\textbf{Step 3.} Handle the non-degenerate case: locate an interior extremum. Assume $M > m$. Since $f(a) = f(b)$, at most one of $\{M, m\}$ can equal the common endpoint value $f(a)$. Therefore at least one of $M, m$ is not equal to $f(a) = f(b)$. Call it $E$ (either the maximum or minimum). Any point at which $f$ attains the value $E$ cannot be an endpoint (since $E \neq f(a) = f(b)$), so it must lie in the open interval $(a,b)$. Let $c \in (a,b)$ be a point where $f(c) = E$.

\textbf{Step 4.} Identify $c$ as a relative extremum. If $E = M$ then $f(c) = M \geq f(x)$ for all $x \in [a,b]$, so in particular $f(c) \geq f(x)$ for all $x$ in any neighbourhood of $c$ within $[a,b]$. Thus $f$ has a relative maximum at $c$. If $E = m$ then $f(c) = m \leq f(x)$ for all $x \in [a,b]$, so $f$ has a relative minimum at $c$. In either case $c$ is a relative extremum of $f$.

\textbf{Step 5.} Apply the Interior Extremum Theorem. Since $c \in (a,b)$, $c$ is an interior point of $[a,b]$. Since $f$ is differentiable on $(a,b)$ and $c \in (a,b)$, $f$ is differentiable at $c$. By the Interior Extremum Theorem, differentiability at a relative extremum implies $f'(c) = 0$.
\end{proof}

\begin{remark*}[Proof structure]
The proof is a case split on whether $f$ is constant, followed by an existence argument for an interior extremum. The constant case is immediate. The non-constant case uses a pigeonhole observation: with $f(a) = f(b)$, both the maximum and minimum cannot simultaneously equal the endpoint value, so at least one extreme value is attained strictly inside the interval. The Interior Extremum Theorem then converts the extremum into a zero of the derivative. The three ingredients are the Extreme Value Theorem (existence of extreme values), a boundary exclusion argument (the extreme value is not at the endpoints), and the Interior Extremum Theorem (differentiability at an interior extremum forces zero derivative).
\end{remark*}

\begin{remark*}[Dependencies]~\\
\begin{itemize}
  \item \hyperref[thm:extreme-value-theorem]{Theorem: Extreme Value Theorem}
  \item \hyperref[thm:interior-extremum-theorem]{Theorem: Interior Extremum Theorem}
  \item \hyperref[def:derivative-at-a-point]{Definition: Derivative at a Point}
\end{itemize}
\end{remark*}
